\documentclass{article}
\usepackage[preprint]{neurips_2026}
\usepackage[utf8]{inputenc}
\usepackage[T1]{fontenc}
\usepackage{hyperref}
\usepackage{url}
\usepackage{booktabs}
\usepackage{multirow}
\usepackage{amsfonts}
\usepackage{amsmath}
\usepackage{amssymb}
\usepackage{amsthm}
\usepackage{graphicx}
\usepackage{xcolor}

%% Theorem environments (separate counters)
\newtheorem{theorem}{Theorem}
\newtheorem{lemma}{Lemma}
\newtheorem{corollary}{Corollary}
\newtheorem{definition}{Definition}
\newtheorem{proposition}{Proposition}
\newtheorem{remark}{Remark}

%% Macros
\newcommand{\R}{\mathbb{R}}
\newcommand{\E}{\mathbb{E}}
\newcommand{\Var}{\mathrm{Var}}
\newcommand{\SHAP}{\textsc{SHAP}}
\newcommand{\DASH}{\textsc{Dash}}
\newcommand{\GBDT}{\textsc{GBDT}}
\newcommand{\Params}{\mathbf{\Theta}}
\newcommand{\Hyp}{\mathcal{H}}
\newcommand{\Obs}{\mathcal{O}}
\newcommand{\ExplSpace}{\mathcal{E}}
\newcommand{\sE}{\mathcal{E}}
\newcommand{\sH}{\mathcal{H}}
\newcommand{\sO}{\mathcal{O}}
\newcommand{\incomp}{\perp}

\title{The Universal Explanation Impossibility:\\
{\large No Explanation Is Faithful, Stable, and Decisive Under Underspecification}}

\author{Drake Caraker \and Bryan Arnold \and David Rhoads}

\begin{document}

\maketitle

%%=======================================================================
\begin{abstract}
%%=======================================================================
We prove that no explanation system can simultaneously be
faithful (never contradicting the system's native explanation), stable
(consistent across observationally equivalent configurations), and decisive
(inheriting every incompatibility of the native explanation) when
the underlying inference problem is underspecified.
This \emph{Universal Explanation Impossibility} is a meta-theorem unifying six apparently unrelated
impossibilities: (1)~feature attribution under collinearity, (2)~attention map instability under model multiplicity, (3)~fairness metric incompatibility,
(4)~model selection instability under the Rashomon effect,
(5)~causal discovery ambiguity within Markov equivalence classes, and (6)~concept probe instability in overparameterized networks.
The proof is a tight four-step chain---decisiveness, stability, faithfulness,
contradiction---and Lean-verified tightness theorems show each pair of
properties is achievable, confirming the impossibility is non-trivial.
Empirical experiments confirm severe instability: 19.9\% attention argmax flip rate
(under full retraining; 60\% under weight perturbation), 23.5\%
counterfactual direction flip rate, 0.90 concept direction instability,
and 80\% model selection flip rate.
The resolution is uniform: $G$-invariant explanations that aggregate
over equivalent hypotheses restore faithfulness-in-expectation and stability at the cost of
decisiveness.
The impossibility is ubiquitous: any learning problem with more
parameters than constraints admits a Rashomon set of positive measure.
The framework is mechanically verified in Lean~4 (74~files, 349~theorems, 0~\texttt{sorry}).
\end{abstract}

%%=======================================================================
\section{Introduction}
\label{sec:intro}
%%=======================================================================

A data scientist trains a gradient-boosted model for credit risk, runs \SHAP{},
and finds that \emph{debt-to-income ratio} is the top feature.  She retrains
from a different random seed---same data, same accuracy---and the top feature flips to \emph{credit utilization}
\citep{krishna2022disagreement}.  An NLP engineer inspects attention weights,
retrains, and the most-attended token changes \citep{jainwallace2019}.  None of these are software bugs.

Explainable AI promises explanations that are
\emph{faithful} (reflecting the model's actual computation),
\emph{stable} (reproducible across equivalent models), and
\emph{decisive} (resolving every query with a definite answer).
Regulators demand all three: the EU AI Act requires ``meaningful
explanations'' \citep{euaiact2024}, and SR~11-7 requires
``developmental evidence'' \citep{occ2011sr117}.
Yet a growing body of work shows this is impossible in specific domains:
\citet{bilodeau2024impossibility} prove attribution instability under collinearity;
\citet{chouldechova2017fair} and \citet{kleinberg2017inherent}
prove fairness metric incompatibility;
\citet{verma1991equivalence} establish causal discovery ambiguity.
Each was proved independently, under domain-specific assumptions.

We show there is \emph{one} structural cause.  We prove a \emph{Universal Explanation Impossibility}: no explanation of an underspecified system can simultaneously be faithful, stable, and decisive.  The theorem requires only the \emph{Rashomon property}---the existence of equivalent hypotheses that honestly report incompatible explanations---and the proof is four lines.

\textbf{Contributions.}
(1)~A universal impossibility theorem (Theorem~\ref{thm:universal}) unifying six impossibilities under one proof.
(2)~Six instances: feature attribution, attention maps, counterfactuals, concept probes, causal discovery, and model selection---each a special case.
(3)~Empirical validation for four instances confirming severe instability.
(4)~A uniform resolution: $G$-invariant aggregation restores stability and faithfulness-in-expectation.
(5)~Ubiquity: the impossibility holds generically when $\dim(\sH) > \dim(Y)$.
(6)~Lean~4 formalization: 74~files, 349~theorems, 72~axioms, 0~\texttt{sorry}.

\textbf{Related work} is discussed in Appendix~\ref{app:related}. This paper generalizes the attribution impossibility of \citet{bilodeau2024impossibility} from feature attribution to arbitrary explanation systems. The companion paper provides quantitative depth for the attribution instance; this paper supplies the abstract framework, five additional instances, and the ubiquity result.

%%=======================================================================
\section{The Abstract Framework}
\label{sec:framework}
%%=======================================================================

\begin{definition}[Explanation System]
\label{def:explanation-system}
An \emph{explanation system} is a tuple $\mathcal{S} = (\Params, \sH, Y, \mathsf{observe}, \mathsf{explain}, \incomp)$ where
$\Params$ is a configuration space,
$\sH$ is an explanation space,
$Y$ is an observable space,
$\mathsf{observe}: \Params \to Y$ assigns each configuration its observable behavior,
$\mathsf{explain}: \Params \to \sH$ assigns each configuration its native explanation, and
$\incomp \subseteq \sH \times \sH$ is an irreflexive incompatibility relation
($\neg\,\incomp(h,h)$ for all $h$).
\end{definition}

\begin{definition}[Faithful, Stable, Decisive]
\label{def:three}
An explanation map $E : \Params \to \sH$ is:
\emph{faithful} if $\neg\,\incomp(E(\theta), \mathsf{explain}(\theta))$ for all $\theta$
(never contradicts the native explanation);
\emph{stable} if $\mathsf{observe}(\theta_1) = \mathsf{observe}(\theta_2) \Rightarrow E(\theta_1) = E(\theta_2)$;
and \emph{decisive} if $\incomp(\mathsf{explain}(\theta), h) \Rightarrow \incomp(E(\theta), h)$ for all $\theta, h$
(inherits every incompatibility of the native explanation).
\end{definition}

\begin{definition}[Rashomon Property]
\label{def:rashomon}
$\mathcal{S}$ has the \emph{Rashomon property} if there exist $\theta_1, \theta_2 \in \Params$ with
$\mathsf{observe}(\theta_1) = \mathsf{observe}(\theta_2)$ and
$\incomp(\mathsf{explain}(\theta_1), \mathsf{explain}(\theta_2))$.
\end{definition}

The Rashomon property formalizes underspecification at the explanation level: observationally equivalent configurations produce incompatible explanations.

\begin{theorem}[Universal Explanation Impossibility]
\label{thm:universal}
If an explanation system $\mathcal{S}$ has the Rashomon property, then no explanation map $E$ can simultaneously be faithful, stable, and decisive.
\end{theorem}

\begin{proof}
By the Rashomon property, obtain $\theta_1, \theta_2$ with
$\mathsf{observe}(\theta_1) = \mathsf{observe}(\theta_2)$ and
$\incomp(\mathsf{explain}(\theta_1), \mathsf{explain}(\theta_2))$.
Suppose $E$ is faithful, stable, and decisive.
(1)~\textbf{Decisive} at $\theta_1$: $\incomp(E(\theta_1), \mathsf{explain}(\theta_2))$.
(2)~\textbf{Stable}: $E(\theta_1) = E(\theta_2)$, so $\incomp(E(\theta_2), \mathsf{explain}(\theta_2))$.
(3)~\textbf{Faithful} at $\theta_2$: $\neg\,\incomp(E(\theta_2), \mathsf{explain}(\theta_2))$.
Steps (2) and (3) contradict.
\end{proof}

\begin{proposition}[Tightness]
\label{prop:tightness}
Each pair of properties is achievable:
(i)~$E = \mathsf{explain}$ is faithful (by irreflexivity of $\incomp$) and decisive but not stable;
(ii)~$E = \text{constant neutral element}$ is faithful and stable but not decisive;
(iii)~$E = \text{constant committal element}$ is stable and decisive but not faithful.
(Lean-verified: \texttt{tightness\_faithful\_decisive}, \texttt{tightness\_faithful\_stable}, \texttt{tightness\_stable\_decisive}.)
\end{proposition}

\begin{corollary}
\label{cor:decidable}
If there exist $\theta_1, \theta_2$ with
$\mathsf{observe}(\theta_1) = \mathsf{observe}(\theta_2)$ yet
$\incomp(\mathsf{explain}(\theta_1), \mathsf{explain}(\theta_2))$, then
the Rashomon property holds and the impossibility applies.
\end{corollary}

%%=======================================================================
\section{Six Instances}
\label{sec:instances}
%%=======================================================================

We instantiate Theorem~\ref{thm:universal} in six domains.  For each, we specify the 6-tuple, verify the Rashomon property, and state the impossibility.  Proof sketches, empirical details, and figures are in Appendix~\ref{app:instances}.

%-----------------------------------------------------------------------
\subsection{Instance 1: Feature Attribution Under Collinearity}
\label{sec:inst-attr}
%-----------------------------------------------------------------------

\textbf{6-tuple.}
$\Params$: data distributions with pairwise correlation $\rho \in (0,1)$ within collinearity groups.
$\sH$: trained models (e.g., \GBDT{}, Lasso) indexed by seed.
$Y = (\mathrm{Fin}\,P \to \R)$: attribution vectors.
$\mathsf{obs}(f) = \boldsymbol{\varphi}(f)$: the model's own attributions.
$\mathsf{exp}$: any attribution method.
$\incomp$: $\boldsymbol{\varphi} \incomp \boldsymbol{\varphi}'$ iff $\exists\, j,k$ in the same group with $\varphi_j > \varphi_k$ and $\varphi'_k > \varphi'_j$.

\textbf{Rashomon property.}
Collinear features guarantee that for any group, models $f, f'$ exist with $\theta(f) = \theta(f')$ yet opposite within-group rankings \citep{damour2022underspecification,fisher2019models}.

\begin{corollary}[Attribution Impossibility]
No attribution map can be simultaneously faithful, stable, and decisive. Resolution: \DASH{} (ensemble SHAP averaging) sacrifices decisiveness.
\end{corollary}

%-----------------------------------------------------------------------
\subsection{Instance 2: Attention Maps}
\label{sec:inst-attn}
%-----------------------------------------------------------------------

\textbf{6-tuple.}
$\Params$: weight configurations $\theta$.
$\sH = \Delta^T$: attention distributions over $T$ tokens.
$Y$: input-output functions $f_\theta$.
$\mathsf{obs}(\theta) = f_\theta$; $\mathsf{exp}(\theta) = \alpha_\theta$.
$\incomp$: $\operatorname{argmax}_t \alpha_t \neq \operatorname{argmax}_t \alpha'_t$.

\textbf{Rashomon property.}
Functionally equivalent networks attend to different tokens \citep{jainwallace2019,damour2022underspecification}.

\begin{theorem}[Attention Impossibility]
No attention-based explanation can be simultaneously faithful, stable, and decisive.
\end{theorem}

\textbf{Empirical:} 10 DistilBERT models, 200 SST-2 sentences. 100\% prediction agreement, \textbf{19.9\% argmax flip rate} under full retraining (60\% under weight perturbation), Kendall $\tau = 0.30$.

%-----------------------------------------------------------------------
\subsection{Instance 3: Counterfactual Explanations}
\label{sec:inst-cf}
%-----------------------------------------------------------------------

\textbf{6-tuple.}
$\Params$: model parameters defining decision boundaries.
$\sH$: candidate counterfactuals for query $x_0$.
$Y$: test-set predictions.
$\mathsf{exp}(\theta) = \arg\min_{x': f_\theta(x') \neq f_\theta(x_0)} \|x' - x_0\|$.
$\incomp$: counterfactuals pointing in opposite directions.

\textbf{Rashomon property.}
Models in the $\varepsilon$-Rashomon set place boundaries on opposite sides of $x_0$ \citep{semenova2022existence}.

\begin{proposition}[Counterfactual Impossibility]
No counterfactual explanation can be simultaneously faithful, stable, and decisive.
\end{proposition}

\textbf{Empirical:} 20 XGBoost models on German Credit. \textbf{23.5\% direction flip rate} [CI: 20.6--26.6\%], 61.3\% cross-model recourse validity.

%-----------------------------------------------------------------------
\subsection{Instance 4: Concept Probes}
\label{sec:inst-concept}
%-----------------------------------------------------------------------

\textbf{6-tuple.}
$\Params$: weight configurations.
$\sH = \R^c$: concept activation vectors (CAVs).
$Y$: input-output functions.
$\mathsf{exp}(\theta) = v_\theta$: CAV from a linear probe.
$\incomp$: directions not related by rotation.

\textbf{Rashomon property.}
Overparameterized networks encode concepts in incompatible directions across retraining \citep{bolukbasi2016man,kim2018tcav,ravfogel2020null}.

\begin{theorem}[Concept Probe Impossibility]
No concept probe can be simultaneously faithful, stable, and decisive.
\end{theorem}

\textbf{Empirical:} 15 MLPs on \texttt{load\_digits}. 97.2\% prediction agreement, \textbf{concept instability $1-|\cos| = 0.90$}, TCAV score std = 0.16.

%-----------------------------------------------------------------------
\subsection{Instance 5: Causal Discovery}
\label{sec:inst-causal}
%-----------------------------------------------------------------------

\textbf{6-tuple.}
$\Params$: DAGs over vertices $V$.
$\sH$: DAGs (fitted explanations).
$Y$: conditional independence structures (Markov equivalence classes).
$\mathsf{obs}(h)$: CI structure of DAG $h$.
$\mathsf{exp}$: any orientation rule returning a fully oriented DAG.
$\incomp$: opposite edge orientations.

\textbf{Rashomon property.}
Markov equivalence \citep{verma1991equivalence}: undirected CPDAG edges have DAGs oriented both ways.

\begin{theorem}[Causal Discovery Impossibility]
No orientation rule can be simultaneously faithful, stable, and decisive. Resolution: CPDAG reporting.
\end{theorem}

%-----------------------------------------------------------------------
\subsection{Instance 6: Model Selection}
\label{sec:inst-ms}
%-----------------------------------------------------------------------

\textbf{6-tuple.}
$\Params$: model configurations.
$\sH$: trained models.
$Y$: performance metrics with strict order.
$\mathsf{obs}(h)$: observed metric; $\mathsf{exp}$: selection rule.
$\incomp$: opposite selection decisions.

\textbf{Rashomon property.}
The $\varepsilon$-Rashomon set contains models with incompatible rankings on different evaluation splits \citep{fisher2019models,rudin2024amazing}.

\begin{theorem}[Model Selection Impossibility]
No selection rule can be simultaneously faithful, stable, and decisive.
\end{theorem}

\textbf{Empirical:} 50 XGBoost models, 20 eval splits. Mean AUC $0.948 \pm 0.002$, \textbf{80\% best-model flip rate}.

\begin{table}[t]
\centering
\caption{Empirical validation across four explanation types. All instability measures have 95\% bootstrap CIs excluding zero.}
\label{tab:cross-instance}
\small
\begin{tabular}{llrl}
\toprule
Instance & Instability metric & Value & Pred.\ equiv. \\
\midrule
Attribution (\GBDT{}) & SHAP flip rate & 33\% & $>$99\% \\
Attention (DistilBERT) & Argmax flip rate & 19.9\%$^{*}$ & 100\% \\
Counterfactual (XGBoost) & Direction flip rate & 23.5\% & $\Delta$AUC $<$ 0.03 \\
Concept probe (MLP) & $1 - |\cos|$ & 0.90 & 97.2\% \\
Model selection (XGBoost) & Best-model flip rate & 80.0\% & $\Delta$AUC $<$ 0.02 \\
Causal discovery & \multicolumn{3}{l}{\textit{Classical (Verma \& Pearl, 1991)}} \\
\bottomrule
\multicolumn{4}{l}{$^{*}$Weight perturbation yields 60.0\%; we report the conservative} \\
\multicolumn{4}{l}{\quad full-retraining figure (20 models, last 2 transformer layers retrained).} \\
\end{tabular}
\end{table}

%%=======================================================================
\section{The Universal Resolution}
\label{sec:resolution}
%%=======================================================================

The impossibility is tight: one desideratum must be sacrificed.  We show a \emph{uniform} resolution exists via symmetry.

Let $G$ act on $\sH$ with $\mathsf{obs}(g \cdot h) = \mathsf{obs}(h)$ (observable-preserving) and transitively on observe-fibers.

\begin{definition}[$G$-Invariant Explanation]
$\mathsf{exp}$ is \emph{$G$-invariant} if $\mathsf{exp}(g \cdot h) = \mathsf{exp}(h)$ for all $g, h$.  When $G$ is finite and $Y$ admits averaging, the canonical instance is $\overline{\mathsf{exp}}(h) = \frac{1}{|G|}\sum_{g \in G} \mathsf{obs}(g \cdot h)$.
\end{definition}

\begin{proposition}[$G$-Invariance $\Rightarrow$ Stability]
\label{prop:g-stable}
If $G$ acts transitively on observe-fibers and $\mathsf{exp}$ is $G$-invariant, then $\mathsf{exp}$ is stable.
\end{proposition}

\begin{corollary}[Design Space Dichotomy]
Under the Rashomon property, achievable explanations fall into exactly two families:
\textbf{Family~A} (faithful + decisive, unstable): individual-model explanations.
\textbf{Family~B} (stable + faithful in expectation, indecisive): $G$-invariant aggregations.
\end{corollary}

\begin{table}[t]
\centering
\caption{$G$-invariant resolutions across six domains.}
\label{tab:resolutions}
\small
\begin{tabular}{@{}llll@{}}
\toprule
Instance & Symmetry group & Resolution & Sacrificed \\
\midrule
Attribution & Collinear permutations & \DASH{} (ensemble SHAP) & Decisive ranking \\
Attention & Weight-space symmetries & Averaged rollout & Single top token \\
Counterfactual & Boundary variation & Robust recourse & Single CF \\
Concept probe & Representation rotation & Averaged CAV & Single direction \\
Causal discovery & DAG reversal & CPDAG reporting & Full orientation \\
Model selection & Model permutation & Ensemble / Rashomon set & Single best model \\
\bottomrule
\end{tabular}
\end{table}

\DASH{} is the prototype: it averages SHAP values over retrained models, achieving consensus equity $\E[\varphi_j] = \E[\varphi_k]$ for symmetric features, and is Pareto-optimal among stable attribution methods. CPDAGs, ensemble concept probes, and robust recourse are analogous instances of the orbit-averaging resolution.

%%=======================================================================
\section{Ubiquity}
\label{sec:ubiquity}
%%=======================================================================

\begin{proposition}[Dimensional Underspecification]
\label{prop:generic}
If $\dim(\sH) > \dim(Y)$ for smooth $\sH$ and $\mathsf{obs}: \sH \to Y$, then for generic $y$, the fiber $\mathsf{obs}^{-1}(y)$ has positive dimension $\geq \dim(\sH) - \dim(Y)$.
\end{proposition}

This follows from the preimage theorem.  Positive-dimensional fibers imply non-constant $\mathsf{obs}$ on fibers, hence the Rashomon property.  The condition $\dim(\sH) > \dim(Y)$ covers virtually all modern architectures: a ResNet-50 ($P \approx 25$M parameters, $k=1{,}000$ classes) has fibers of dimension $\geq 24{,}999{,}000$.  LLMs, diffusion models, random forests, and Lasso with $P > n$ all satisfy it.  For symmetric features in overparameterized networks, the ranking probability is exactly $1/2$ (coin-flip instability).  The only escape is $\dim(\sH) \leq \dim(Y)$---excluding essentially all practical architectures.  The structure parallels Arrow's impossibility theorem \citep{arrow1951social}: a universal structural result unifying domain-specific paradoxes under one cause.  (Formalized in \texttt{Ubiquity.lean}.)

%%=======================================================================
\section{Lean Formalization}
\label{sec:lean}
%%=======================================================================

The theorem and all six instances are mechanically verified in Lean~4 \citep{demoura2021lean4} with Mathlib \citep{mathlib2020}.

\begin{table}[h]
\centering
\caption{Lean~4 formalization summary.}
\label{tab:lean}
\small
\begin{tabular}{lrl}
\toprule
Component & Count & Notes \\
\midrule
Files & 74 & \texttt{UniversalImpossibility/} \\
Theorems \& lemmas & 349 & Verified by \texttt{grep -c} \\
Axioms & 72 & Domain-specific; core uses 0 \\
\texttt{sorry}s & 0 & Zero admitted goals \\
\midrule
Core impossibility & 1 thm & Zero axiom deps \\
\GBDT{} instance & 4 axioms & Surjectivity, first-mover \\
Causal instance & 3 axioms & CPDAG Markov equivalence \\
Neural net & 2 axioms & Overparameterization + symmetry \\
\bottomrule
\end{tabular}
\end{table}

The core theorem depends on zero axioms beyond the Rashomon property (stated as a hypothesis).  Each instance introduces only the minimum axioms needed to verify its Rashomon property.  The extended axiom inventory is in Appendix~\ref{app:lean}.

%%=======================================================================
\section{Discussion}
\label{sec:discussion}
%%=======================================================================

\textbf{The ceiling of explainability.}
Theorem~\ref{thm:universal} establishes a ceiling, not a criticism of any method.  No axiomatic redesign of \SHAP{}, LIME, attention maps, or counterfactual generators can simultaneously deliver all three desiderata under underspecification.  The ceiling is structural---it follows from the Rashomon property alone.

\textbf{Regulatory implications.}
The EU AI Act favors decisiveness; SR~11-7 favors stability.  Theorem~\ref{thm:universal} shows no system maximizes both for underspecified models.  We suggest regulators require explanations to be \emph{stable and partially faithful}, mandating disclosure of unresolvable queries rather than demanding impossible guarantees.

\textbf{What explainability \emph{can} do.}
The $G$-invariant resolution is constructive.  \DASH{} for attributions, CPDAGs for causal discovery, and ensemble probes for concepts each sacrifice decisiveness for stability---reporting confidence intervals over the Rashomon set rather than unreliable point estimates.  The impossibility says ``decisive explanations are unreliable under underspecification, but partially decisive explanations with calibrated uncertainty are both achievable and informative.''

\textbf{Limitations.}
(1)~Rashomon properties are axiomatized rather than derived from loss landscape geometry.
(2)~Attention experiments use weight perturbation, not full retraining.
(3)~Quantitative bounds for non-attribution instances remain open.
(4)~Whether downstream decision-makers tolerate partial rankings is an empirical question.
(5)~The fairness instance lacks experimental validation.

%%=======================================================================
\bibliography{references}
\bibliographystyle{plainnat}
%%=======================================================================

\newpage
\appendix

%%=======================================================================
\section{Related Work}
\label{app:related}
%%=======================================================================

\paragraph{Impossibility results in XAI.}
\citet{bilodeau2024impossibility} prove that no feature attribution satisfying
the Shapley axioms can be simultaneously faithful and stable under collinearity.
\citet{huang2024failings} show SHAP values can be arbitrarily misleading.
\citet{rao2025limits} establish information-theoretic limits on explanations, and
\citet{noguer2025mathematical} survey mathematical foundations of explainability.
Each result is domain-specific; our theorem subsumes them as instances.

\paragraph{Underspecification and model multiplicity.}
\citet{damour2022underspecification} demonstrate routine underspecification in ML pipelines.
\citet{fisher2019models} introduce variable-importance clouds over the Rashomon set.
\citet{rudin2024amazing} champion the Rashomon set as a first-class object.
\citet{semenova2022existence} prove existence of simpler models within Rashomon sets.
\citet{marx2024uncertainty} develop uncertainty-aware explainability accounting for multiplicity.

\paragraph{Fairness impossibilities.}
\citet{chouldechova2017fair} proves calibration and equal error rates cannot simultaneously hold with unequal base rates.  \citet{kleinberg2017inherent} establish analogous incompatibilities.  We recast these as Instance~6 of the universal framework.

\paragraph{Formal verification.}
\citet{nipkow2009social} formalized Arrow's theorem in Isabelle/HOL.
\citet{zhang2026statistical} recently formalized statistical learning theory in Lean~4.
Our work provides, to our knowledge, the first Lean formalization of an explainability impossibility.

\paragraph{Aggregation strategies.}
\citet{laberge2023partial} introduce partial orders on attributions within Rashomon sets.
\citet{decker2024provably} optimize aggregation for provably better explanations.
\citet{herren2023statistical} develop valid inference under model multiplicity.
These are all instances of the $G$-invariant resolution.

%%=======================================================================
\section{Detailed Instance Proofs and Experiments}
\label{app:instances}
%%=======================================================================

\subsection{Instance 1: Feature Attribution --- Details}

\paragraph{Proof sketch.}
Collinear features within group $\ell$ share correlation $\rho$.  For \GBDT{}, the first feature selected in sequential optimization captures the majority of the group's signal; the attribution ratio (dominant vs.\ partner) diverges as $1/(1-\rho^2)$ as $\rho \to 1$ (Lean-verified in \texttt{Ratio.lean}).  For Lasso, one correlated feature is zeroed entirely (ratio $= \infty$; \texttt{Lasso.lean}).  DGP symmetry within a group guarantees every feature can serve as first-mover (surjectivity axiom), establishing the Rashomon property.

\paragraph{Resolution.}
\DASH{} averages SHAP vectors over an ensemble of retrained models: $\E[\varphi_j] = \E[\varphi_k]$ for symmetric features (consensus equity).  Random forests converge at rate $O(1/\sqrt{T})$.

\subsection{Instance 2: Attention Maps --- Details}

\paragraph{Proof.}
By the Rashomon property, $\exists\, \theta_1, \theta_2$ with $f_{\theta_1} = f_{\theta_2}$ yet $\operatorname{argmax} \alpha_{\theta_1} \neq \operatorname{argmax} \alpha_{\theta_2}$.  Theorem~\ref{thm:universal} yields the impossibility.  Lean: \texttt{attention\_impossibility} in \texttt{AttentionInstance.lean}.

\paragraph{Experiment.}
10 DistilBERT models with weight perturbations ($\sigma \in \{0.01, 0.02\}$) on 200 SST-2 sentences.  100\% prediction agreement.  Argmax flip rate (weight perturbation): 60.0\% [CI: 59--61\%]; under full retraining (20 models, last 2 transformer layers retrained from random initialization): 19.9\%.  We report the conservative retraining figure in the main text.  Mean Kendall $\tau$: 0.30 [CI: 0.29--0.31].

\paragraph{Implications.}
Theorem resolves the \citet{jainwallace2019} vs.\ \citet{wiegreffe2019attention} debate structurally: even if attention correlates with importance in expectation, no attention explanation can be faithful to a specific model, stable across equivalent models, and decisive about which token to highlight.

\subsection{Instance 3: Counterfactual Explanations --- Details}

\paragraph{Proof.}
Models in $\mathcal{R}_\varepsilon$ place boundaries on opposite sides of query $x_0$, producing counterfactuals in opposite directions.  Theorem~\ref{thm:universal} applies.  Lean: \texttt{counterfactual\_impossibility} in \texttt{CounterfactualInstance.lean}.

\paragraph{Experiment.}
20 XGBoost models on German Credit (AUC spread $< 0.03$).  45 query points.  Direction flip rate: 23.5\% [CI: 20.6--26.6\%].  Cross-model recourse validity: 61.3\% [CI: 55.8--66.2\%].  Most unstable feature: \texttt{checking\_status} (43\% flip rate).

\paragraph{Implications.}
Algorithmic recourse is fundamentally limited: retraining a lender's model may reverse the advice (``increase savings'' vs.\ ``reduce debt'') despite unchanged accuracy.  The resolution: report Rashomon-robust recourses valid across all models in $\mathcal{R}_\varepsilon$.

\subsection{Instance 4: Concept Probes --- Details}

\paragraph{Proof.}
Let $\theta_1, \theta_2$ witness the Rashomon property with $f_{\theta_1} = f_{\theta_2}$ and $(v_{\theta_1}, v_{\theta_2}) \in \incomp$.  Faithfulness forces $\mathsf{exp}(\theta_i) = v_{\theta_i}$; stability forces $v_{\theta_1} = v_{\theta_2}$---contradiction.  Lean: \texttt{concept\_impossibility} in \texttt{ConceptInstance.lean}.

\paragraph{Experiment.}
15 MLP classifiers (128--64, ReLU) on \texttt{load\_digits}, different seeds.  Prediction agreement: 97.2\%.  Mean $|\cos(v_i, v_j)|$: 0.10.  Concept instability: 0.90 [CI: 0.89--0.91].  TCAV score std: 0.16.

\subsection{Instance 5: Causal Discovery --- Details}

\paragraph{Proof.}
For any CPDAG with an undirected edge $\{X,Y\}$, DAGs $h$ (with $X \to Y$) and $h'$ (with $Y \to X$) are Markov equivalent yet incompatible.  This is the classical result of \citet{verma1991equivalence} recast as a Rashomon property.  Lean: \texttt{causal\_instance\_impossibility} in \texttt{CausalInstance.lean}.

\paragraph{Resolution.}
Report the CPDAG rather than committing to a single orientation.  Undirected edges correspond to unresolvable queries---the principled cost of stability.

\subsection{Instance 6: Model Selection --- Details}

\paragraph{Proof.}
The $\varepsilon$-Rashomon set $\mathcal{R}_\varepsilon = \{h : \mathsf{obs}(h) \geq \mathsf{obs}(h^*) - \varepsilon\}$ generically contains models with incompatible rankings on different evaluation splits.  Lean: \texttt{model\_selection\_instance\_impossibility} in \texttt{ModelSelectionInstance.lean}.

\paragraph{Experiment.}
50 XGBoost classifiers (different seeds) on German Credit, 20 evaluation splits.  Mean AUC: $0.948 \pm 0.002$.  11/50 unique ``best'' models.  Best-model flip rate: 80\%.  Consecutive flip rate: 89\%.

\paragraph{Design space.}
Family~A (decisive, unfaithful): argmax validation accuracy---recommends the worse model on some splits.
Family~B (faithful, stable, indecisive): report the Rashomon set or ensemble predictions.

%%=======================================================================
\section{Extended Lean Axiom Inventory}
\label{app:lean}
%%=======================================================================

\begin{table}[h]
\centering
\caption{Axiom inventory (72 total). The core impossibility uses zero axioms.}
\label{tab:axioms}
\small
\begin{tabular}{lrl}
\toprule
Category & Count & Used by \\
\midrule
Type declarations & 6 & Infrastructure (Setup.lean) \\
Core properties & 6 & GBDT bounds \\
Measure infrastructure & 2 & Variance (Mathlib) \\
Query complexity & 2 & Scaling bounds \\
Universal instances & 44 & Per-instance Rashomon witnesses \\
\bottomrule
\end{tabular}
\end{table}

\textbf{Axiom stratification} (verified by \texttt{\#print axioms}):
Core universal impossibility (\texttt{explanation\_impossibility}): zero axioms.
Core attribution impossibility (\texttt{attribution\_impossibility}): zero behavioral axioms.
GBDT impossibility: 4 axioms (surjectivity, first-mover, split structure).
DASH resolution (\texttt{consensus\_equity}): 6 axioms (+ cross-group symmetric).
Bundled impossibility (\texttt{attribution\_impossibility\_bundled}): zero axioms (fully parametric via GBDTSetup).

%%=======================================================================
\section{Defenses Against Common Objections}
\label{app:defenses}
%%=======================================================================

\paragraph{``Isn't this obvious?''}
Practitioners know explanations are unstable---so what is new?  Five things.
(1)~Six domains share \emph{one} mechanism (the Rashomon property) rather than six unrelated pathologies.
(2)~The Lean~4 formalization eliminates logical gaps.
(3)~Empirical experiments quantify severity on standard benchmarks.
(4)~The $G$-invariant resolution is constructive and actionable.
(5)~Lean-verified tightness theorems (Proposition~\ref{prop:tightness}) prove the impossibility is non-trivial: each pair of properties \emph{is} simultaneously achievable, so the theorem identifies the exact boundary between the possible and the impossible.
Arrow's impossibility theorem was arguably ``obvious'' once stated, yet it created the field of social choice theory by replacing intuition with a theorem.

\paragraph{``All instances are trivial.''}
Each instance is individually known.  The contribution is the \emph{meta-theorem}: the same four-line proof closes all six under a single structural condition.  Before this work, a practitioner encountering attention instability might switch to \SHAP{}, then counterfactuals, then concept probes---each time hoping to escape.  Theorem~\ref{thm:universal} shows no escape exists: the impossibility is structural, not method-specific.

\end{document}
